% ============================================================
\chapter{Markov-Kakutani Theorem}
\label{ch:markov_kakutani}
% ============================================================

\begin{intuition}
The Markov-Kakutani theorem guarantees the existence of a common fixed point for a family
of commuting continuous affine maps acting on a compact convex set in a locally convex
topological vector space. This result, simple in its statement, has deep applications: it
underlies the theory of invariant means and plays a central role in the study of amenable
groups.
\end{intuition}

% ------------------------------------------------------------
\section{Statement and Proof of the Theorem}
% ------------------------------------------------------------

\begin{definition}[Affine map]
\index{affine map}
Let $E$ be a vector space and $C \subseteq E$ a convex set. A map $T\colon C \to C$ is
\emph{affine} if for all $x, y \in C$ and all $\lambda \in [0,1]$:
\[
  T(\lambda x + (1-\lambda)y) = \lambda T(x) + (1-\lambda) T(y).
\]
\end{definition}

\begin{theorem}[Markov-Kakutani]
\label{thm:markov-kakutani}
\index{Markov-Kakutani theorem}
Let $E$ be a locally convex topological vector space and $C \subseteq E$ a nonempty compact
convex set. Let $\mathcal{F}$ be a family of continuous affine maps $T\colon C \to C$ that
pairwise commute ($T_1 \circ T_2 = T_2 \circ T_1$ for all $T_1, T_2 \in \mathcal{F}$).
Then there exists $x^* \in C$ such that $T(x^*) = x^*$ for all $T \in \mathcal{F}$.
\end{theorem}

\begin{proof}
\emph{Step 1: single map case.}
Let $T\colon C \to C$ be continuous and affine. For each $x \in C$ and $n \geq 1$, set
\[
  \sigma_n(x) = \frac{1}{n} \sum_{k=0}^{n-1} T^k(x).
\]
Since $C$ is convex and $T(C) \subseteq C$, we have $\sigma_n(x) \in C$. Define
$C_n = \sigma_n(C) = \{\sigma_n(x) : x \in C\}$. Then $C_n$ is compact (continuous image
of a compact set) and nonempty.

We show that $(C_n)_{n \geq 1}$ has the finite intersection property, which implies
$\bigcap_n \overline{C_n} \neq \emptyset$ by compactness of $C$.

Let $y \in \bigcap_n \overline{C_n}$. For every $\varepsilon$-neighborhood $V$ of $0$,
there exists $x_n \in C$ with $\sigma_n(x_n) - y \in V$. One verifies that
$T(\sigma_n(x_n)) - \sigma_n(x_n) = \frac{1}{n}(T^n(x_n) - x_n) \to 0$ since $C$ is
bounded. Thus $T(y) = y$.

\emph{Step 2: commuting family.}
For each $T \in \mathcal{F}$, define $\mathrm{Fix}(T) = \{x \in C : T(x) = x\}$. By
Step~1, $\mathrm{Fix}(T)$ is nonempty. Moreover:
\begin{itemize}
  \item $\mathrm{Fix}(T)$ is closed (preimage of the closed diagonal under $(x, T(x))$)
    and contained in the compact $C$, hence compact.
  \item $\mathrm{Fix}(T)$ is convex since $T$ is affine: if $T(x) = x$ and $T(y) = y$,
    then $T(\lambda x + (1-\lambda)y) = \lambda x + (1-\lambda)y$.
\end{itemize}
If $T_1, T_2$ commute and $x \in \mathrm{Fix}(T_1)$, then
$T_1(T_2(x)) = T_2(T_1(x)) = T_2(x)$, so
$T_2(\mathrm{Fix}(T_1)) \subseteq \mathrm{Fix}(T_1)$. Thus $T_2$ restricted to
$\mathrm{Fix}(T_1)$ is a continuous affine map of a compact convex set into itself, and
by Step~1 it has a fixed point in $\mathrm{Fix}(T_1)$.

By transfinite induction (or the finite intersection property),
$\bigcap_{T \in \mathcal{F}} \mathrm{Fix}(T) \neq \emptyset$.
\end{proof}

% ------------------------------------------------------------
\section{Amenable Groups}
% ------------------------------------------------------------

\begin{definition}[Invariant mean]
\index{invariant mean}
Let $G$ be a discrete group. A \emph{left-invariant mean} on $G$ is a linear functional
$m\colon \ell^\infty(G) \to \R$ such that:
\begin{enumerate}
  \item $m(\ind_G) = 1$ (normalization);
  \item $m(f) \geq 0$ if $f \geq 0$ (positivity);
  \item $m(L_g f) = m(f)$ for all $g \in G$ (invariance), where $(L_g f)(x) = f(g^{-1}x)$.
\end{enumerate}
\end{definition}

\begin{definition}[Amenable group]
\index{amenable group}
A discrete group $G$ is \emph{amenable} if it admits a left-invariant mean.
\end{definition}

\begin{theorem}[Abelian groups are amenable]
\label{thm:abelian-amenable}
Every discrete abelian group is amenable.
\end{theorem}

\begin{proof}
Let $G$ be an abelian group. The set of means
\[
  \mathcal{M} = \{m \in (\ell^\infty(G))^* : m \geq 0, \; m(\ind_G) = 1\}
\]
is a nonempty compact convex set for the weak-$*$ topology (by the Banach--Alaoglu theorem).

For each $g \in G$, the map $T_g\colon \mathcal{M} \to \mathcal{M}$ defined by
$(T_g m)(f) = m(L_{g^{-1}} f)$ is affine and continuous (for the weak-$*$ topology). It
maps $\mathcal{M}$ into $\mathcal{M}$ since $L_{g^{-1}}$ preserves positivity and
$L_{g^{-1}} \ind_G = \ind_G$.

Since $G$ is abelian, $L_g \circ L_h = L_{gh} = L_{hg} = L_h \circ L_g$, so the $T_g$
commute. By the Markov-Kakutani theorem, there exists
$m^* \in \bigcap_{g \in G} \mathrm{Fix}(T_g)$, i.e., an invariant mean.
\end{proof}

\begin{example}[Invariant mean on $\Z$]
The Banach mean on $\Z$ constructed above is not explicit. However, one can show that for
any bounded sequence $(a_n)_{n \in \Z}$, if the Ces\`aro limit
$\lim_{N \to \infty} \frac{1}{2N+1} \sum_{n=-N}^{N} a_n$ exists, then it coincides with
$m^*((a_n))$.
\end{example}

% ------------------------------------------------------------
\section{Examples of Non-Amenable Groups}
% ------------------------------------------------------------

\begin{theorem}[von Neumann]
The free group on two generators $F_2$ is not amenable.
\end{theorem}

\begin{proof}[Proof sketch]
One uses the Banach--Tarski paradox: if $F_2$ were amenable, one could construct a finitely
additive invariant measure on $F_2$, contradicting the fact that $F_2$ admits a paradoxical
decomposition (a partition into four pieces that, by translations, reconstitute two copies
of $F_2$).
\end{proof}

\begin{proposition}[Properties of amenability]
\begin{enumerate}
  \item Subgroups of amenable groups are amenable.
  \item Quotients of amenable groups are amenable.
  \item Extensions of amenable groups by amenable groups are amenable.
  \item Every finite group is amenable (take the uniform mean).
  \item A group containing $F_2$ as a subgroup is not amenable.
\end{enumerate}
\end{proposition}

% ------------------------------------------------------------
\section{Variants and Generalizations}
% ------------------------------------------------------------

\begin{theorem}[Ryll-Nardzewski]
\label{thm:ryll-nardzewski}
\index{Ryll-Nardzewski theorem}
Let $E$ be a Banach space and $C \subseteq E$ a nonempty weakly compact convex set. Let
$\mathcal{F}$ be a semigroup of affine isometries of $C$. Then
$\bigcap_{T \in \mathcal{F}} \mathrm{Fix}(T) \neq \emptyset$.
\end{theorem}

\begin{remark}
The Ryll-Nardzewski theorem generalizes Markov-Kakutani by replacing commutativity with a
nonexpansiveness condition and weakening compactness to weak compactness.
\end{remark}

\begin{keyformulas}[title=\textbf{Common fixed points}]
\begin{center}
\renewcommand{\arraystretch}{1.3}
\begin{tabular}{|l|l|l|}
\hline
\textbf{Theorem} & \textbf{Family} & \textbf{Space} \\
\hline
Markov-Kakutani & Commuting affine maps & Compact convex, lctvs \\
Ryll-Nardzewski & Affine isometries & Weakly compact convex, Banach \\
\hline
\end{tabular}
\end{center}
\end{keyformulas}

% ------------------------------------------------------------
\section{Exercises}
% ------------------------------------------------------------

\begin{exercise}
Show that the Markov-Kakutani theorem is false without the commutativity hypothesis.
\emph{Hint: consider two rotations of the disk.}
\end{exercise}

\begin{exercise}
Apply the Markov-Kakutani theorem to show that every compact abelian group $G$ admits a
Haar measure (by taking $C$ to be the set of probability measures on $G$).
\end{exercise}

\begin{exercise}
Show that the isometry group of the Euclidean plane is not amenable.
\emph{Hint: show that it contains $F_2$ as a subgroup.}
\end{exercise}

\begin{exercise}
Let $T\colon C \to C$ be a continuous affine map on a compact convex set $C$ in a Banach
space. Show directly (without Markov-Kakutani) that the Ces\`aro means
$\sigma_n(x) = \frac{1}{n}\sum_{k=0}^{n-1} T^k(x)$ converge (at least along a
subsequence) to a fixed point of $T$.
\end{exercise}

\begin{exercise}
Let $G$ be a discrete amenable group acting on a compact space $X$ by homeomorphisms.
Show that there exists a Borel probability measure on $X$ invariant under the action of $G$.
\end{exercise}

\begin{exercise}[Amenability and F\o lner sequences]
\index{F\o lner sequence}
Show that a countable group $G$ is amenable if and only if it admits a \emph{F\o lner
sequence}: a sequence $(F_n)$ of nonempty finite subsets of $G$ such that for every
$g \in G$,
\[
  \frac{\abs{gF_n \triangle F_n}}{\abs{F_n}} \to 0 \quad \text{as } n \to \infty.
\]
Verify that $F_n = \{-n, \ldots, n\}$ is a F\o lner sequence for $\Z$.
\end{exercise}
